\section{Discussion}
\label{sec:discussion}

The main finding is that a RoPE base does not exhaust the design of a finite
frequency table. Support $(a,R)$ determines the sampled interval; allocation
$z$ determines how many channels are spent in each part of it. A scalar-base
sweep explores only the geometric path $z_k=k/(K{-}1)$, whereas the
fixed-support intervention moves outside that path without changing either
endpoint. The resulting change in trained behaviour makes allocation a real
design variable rather than an alternative notation for base selection.

\paragraph{Why the discovery matters.}
The spectral-budget identity reveals the pressure behind this variable. Slow
geometric pairs collapse onto nearly the same positional subspace, so nominal
rotary dimension can substantially overstate effective positional dimension.
The effect becomes most visible when channels are scarce: with $K{=}16$, the
three-seed $432$M MLA comparison lowers $16$K PPL in every seed. A closed-form
allocation makes this finite-budget decision explicit and reproducible, while
learned-frequency methods combine the table choice with an additional
optimization problem.

The construction remains consequential across realistic model-building paths.
At $454$M, the same \rs{} operator has greater leverage on the \evq{} substrate;
at $750$M and $1.485$B, Geo leads in-window before \evq{} overtakes it at longer
lengths. Frozen interventions and low-rank adaptation extend the result to
pretrained checkpoints. These protocols supply architecture, training-stage,
and scale evidence around the fixed-support identification. At $8$B, blocking
answer-side attention to the distant source while preserving prompt tokens and
rotary indices adds a causal capability endpoint: the improved long-position
probability depends on that source.

Static geometry and learned coefficients jointly determine how a frequency
table behaves with fixed weights. The $50$M crossing makes the interaction
visible: a table with better static rank can still destroy perplexity.
Theorem~\ref{thm:obstruction} characterises its exact limit---fixed invertible
Q/K maps preserve all relative-position logits only when the frequency
multisets agree. The frozen OLMo/Qwen controls address the practical
finite-distribution case: changing $z$ at fixed support and amplitude remains
consequential after pretraining.

Support selection and range transport remain complementary. The former chooses
the sampled spectrum, allocation chooses its finite internal resolution, and a
range operator redeploys the learned spectrum at a target length. The common
\rs{} operator has greater leverage on the \evq{} substrate, while the
target-retargeted control in App.~\ref{sec:identification-details} confirms that
support and allocation interact when chosen jointly. LeRoPE supplies compatible
evidence from optimization: a table learned in one run retains substantial
value when frozen for a new training run \citep{karypis2026lerope}. Here, the
closed-form construction removes that table-search stage and the fixed-support
experiment identifies what the table itself changes.

RoPE has a spectral budget because every head has finitely many rotary pairs.
Geometric spacing is one way to spend that budget, not a neutral consequence of
the rotary operator. Once support and allocation are separated, the latter can
be analyzed, controlled, and improved as a first-class part of long-context
training.
